% Options for packages loaded elsewhere
\PassOptionsToPackage{unicode}{hyperref}
\PassOptionsToPackage{hyphens}{url}
%
\documentclass[
]{book}
\usepackage{amsmath,amssymb}
\usepackage{iftex}
\ifPDFTeX
  \usepackage[T1]{fontenc}
  \usepackage[utf8]{inputenc}
  \usepackage{textcomp} % provide euro and other symbols
\else % if luatex or xetex
  \usepackage{unicode-math} % this also loads fontspec
  \defaultfontfeatures{Scale=MatchLowercase}
  \defaultfontfeatures[\rmfamily]{Ligatures=TeX,Scale=1}
\fi
\usepackage{lmodern}
\ifPDFTeX\else
  % xetex/luatex font selection
\fi
% Use upquote if available, for straight quotes in verbatim environments
\IfFileExists{upquote.sty}{\usepackage{upquote}}{}
\IfFileExists{microtype.sty}{% use microtype if available
  \usepackage[]{microtype}
  \UseMicrotypeSet[protrusion]{basicmath} % disable protrusion for tt fonts
}{}
\makeatletter
\@ifundefined{KOMAClassName}{% if non-KOMA class
  \IfFileExists{parskip.sty}{%
    \usepackage{parskip}
  }{% else
    \setlength{\parindent}{0pt}
    \setlength{\parskip}{6pt plus 2pt minus 1pt}}
}{% if KOMA class
  \KOMAoptions{parskip=half}}
\makeatother
\usepackage{xcolor}
\usepackage{longtable,booktabs,array}
\usepackage{calc} % for calculating minipage widths
% Correct order of tables after \paragraph or \subparagraph
\usepackage{etoolbox}
\makeatletter
\patchcmd\longtable{\par}{\if@noskipsec\mbox{}\fi\par}{}{}
\makeatother
% Allow footnotes in longtable head/foot
\IfFileExists{footnotehyper.sty}{\usepackage{footnotehyper}}{\usepackage{footnote}}
\makesavenoteenv{longtable}
\usepackage{graphicx}
\makeatletter
\def\maxwidth{\ifdim\Gin@nat@width>\linewidth\linewidth\else\Gin@nat@width\fi}
\def\maxheight{\ifdim\Gin@nat@height>\textheight\textheight\else\Gin@nat@height\fi}
\makeatother
% Scale images if necessary, so that they will not overflow the page
% margins by default, and it is still possible to overwrite the defaults
% using explicit options in \includegraphics[width, height, ...]{}
\setkeys{Gin}{width=\maxwidth,height=\maxheight,keepaspectratio}
% Set default figure placement to htbp
\makeatletter
\def\fps@figure{htbp}
\makeatother
\setlength{\emergencystretch}{3em} % prevent overfull lines
\providecommand{\tightlist}{%
  \setlength{\itemsep}{0pt}\setlength{\parskip}{0pt}}
\setcounter{secnumdepth}{5}
\usepackage{booktabs}
\ifLuaTeX
  \usepackage{selnolig}  % disable illegal ligatures
\fi
\usepackage[]{natbib}
\bibliographystyle{plainnat}
\usepackage{bookmark}
\IfFileExists{xurl.sty}{\usepackage{xurl}}{} % add URL line breaks if available
\urlstyle{same}
\hypersetup{
  pdftitle={Skilled Reflection},
  hidelinks,
  pdfcreator={LaTeX via pandoc}}

\title{Skilled Reflection}
\author{}
\date{\vspace{-2.5em}}

\begin{document}
\maketitle

{
\setcounter{tocdepth}{1}
\tableofcontents
}
\begin{center}\includegraphics[width=4.17in]{figs/sr_gray_transp} \end{center}

\chapter*{Introduction}\label{introduction}
\addcontentsline{toc}{chapter}{Introduction}

\section{You}\label{you}

You and I:

\begin{verbatim}
0.  **Introduction.** This book is centered on you, the specific reader that you are aware is currently reading. I am the one writing.
1.  **SELF.** For now, focus on applying these words to your actual [SELF][SELF], and determining exactly what or who that is.
2.  **PRI.** Identify and own your [PRIs][Priorities] before you die.
3.  **PPL.** You are not all the other imagined readers. Perhaps they are [PPL][PPL]. You will be too, but not yet.
4.  **BET.** First you'll [BET][BET] about who you (think you) are, until you get sick of losing.
5.  **Education.** At that point you should become [a STUDENT][Education] of better methods.
6.  **WORDs.** As a STUDENT you'll simulate methods efficiently with [WORDs].
7.  **Revision.** You want an answer to the same questions (below) as the rest. [Revise][Revisions] your WORDs until they are uniquely you.
8.  **COMM.** Then you may contextualize and [communicate][Communication] who you are among PPL.  
\end{verbatim}

Until then, you and I and they are all the same.

\section{Questions}\label{questions}

Any questions?

\begin{verbatim}
0.  You. How am I and what do I need?
1.  Self. What must I do to have it?
2.  Priorities. What actions maximize my life?
3.  Others. What are people?
4.  Bet. How do I take action?
5.  Learn. How do I figure out a better plan for my life?
6.  Words. What are they good for?
7.  Revision. What would my life look like on paper?
8.  Communication. How do I relate to or help others?  
\end{verbatim}

\section{Terms}\label{terms}

A book with defined terms is a book for humans.
The following list describes the critical concepts that this book deals with and by impliciation, what it does not deal with.
After each definition is a link to its primary context in the book.
Note that there may be slight discrepancies with definitions in main text as sections become updated. (Last revised: v2.15).

\begin{itemize}
\tightlist
\item
  A GOAL names an ideal future outcome, and requires a PLAN to intentionally achieve it. ({[}\textbf{c2.38}{]}{[}Core Concepts{]})\\
\item
  A PLAN comprises WORDS that facilitate your GOAL. ({[}\textbf{c2.44}{]}{[}Core Concepts{]})\\
\item
  A COST is what's materially required for GOALs--an amount and kind of FORCE. ({[}\textbf{c2.49}{]}{[}Core Concepts{]})\\
\item
  FOOD is the necessary and sufficient material to satisfy these needs. (\hyperref[self]{\textbf{c2.56}})\\
\item
  CAKE is all non-FOOD GOALs. (\hyperref[self]{\textbf{c2.57}})\\
\item
  SELF is the GOAL to MAINTAIN HABIT, ideally to maximize BET outcomes. (\hyperref[self]{\textbf{c2.58}})\\
\item
  HOME is the collection of tools, materials and money (which support HABITs). (\hyperref[home]{\textbf{c2.59}})\\
\item
  PRI is management of the GOALS in your life (and common themes therein). ({[}\textbf{c3.33}{]}{[}Attention and time{]})\\
\item
  PPL are different VERSIONs of each others' PRI (WRITER and READER are different, too).({[}\textbf{c4.1}{]}{[}PPL{]})\\
\item
  TIME is the constant (or denominator) for PRI and FORCES (LIB and MEMORY). ({[}\textbf{c3.38}{]}{[}Attention and time{]})\\
\item
  MONEY is traded for maintenance of, or insurance for SELF and CAKE. ({[}\textbf{c4.09}{]}{[}PPL{]})\\
\item
  WORK is performing a specific task on a collective garden for MONEY. (\hyperref[work]{\textbf{c4.10}})\\
\item
  NORMS are the ``average'' person's HABITs, including WORK and CHUD. (\hyperref[work]{\textbf{c4.13}})\\
\item
  A FORCE is any cause of change including actions. (\hyperref[work]{\textbf{c4.20}})\\
\item
  MAINTENANCE is any recurring COST to neither move toward nor away from a GOAL. (\hyperref[work]{\textbf{c4.21}})\\
\item
  An ALIGNMENT is a reduced COST due to a FORCE shared between at least two GOALs. (\hyperref[work]{\textbf{c4.22}})\\
\item
  A relationship (RLTP) is a GOAL about PPL. (\hyperref[work]{\textbf{c4.23}})\\
\item
  A BET is a PLAN that leverages TIME against CHUD to facilitate GOALs or a REALITY check. (\hyperref[bet]{c5.28})\\
\item
  CHUD is an acronym for assessing the reality of GOALs: COSTs, HABITs, UNKNOWNs and DOUBTs. ({[}\textbf{c5.41}{]}{[}C.H.U.D.{]})\\
\item
  A HABIT is the historical average of your actions and their by-products.({[}\textbf{c5.43}{]}{[}COSTs and HABITs{]})\\
\item
  An UNKNOWN is a FORCE (especially out of your CONTROL) that makes your PLAN wrong.({[}\textbf{c5.44}{]}{[}UNKNOWNs and DOUBTs{]})\\
\item
  A DOUBT is a source of predictable or nameable UNKNOWN, for example, ``denial'' or ``anxiety''. ({[}\textbf{c5.48}{]}{[}UNKNOWNs and DOUBTs{]})\\
\item
  A STUDENT is a WRITER sacrificing TIME for concentrated REVISION of their PRIs. (\hyperref[student]{\textbf{c6.15}})\\
\item
  A READER is a WRITER studying their own reaction to and potential use of others' PLANs. (\hyperref[read]{\textbf{c6.39}})\\
\item
  A WRITER invests TIME and energy to map feelings onto WORDS. ({[}\textbf{c6.64}{]}{[}WRITE PLAN{]})\\
\item
  An INSTRUCTOR advises and enforces STUDENTs' BETs on PLANs (\hyperref[instructor]{\textbf{c6.73}})\\
\item
  A WORD is a BET for simulating (direct attention to) truth (reality) or falsehood (fantasy). ({[}\textbf{c7.20}{]}{[}WORD as BET{]})\\
\item
  An EXAMPLE is an individual, particular event or object of reality. ({[}\textbf{c7.45}{]}{[}EXAMPLE and IDEA{]})\\
\item
  A USE-CASE is an EXAMPLE with consequence (vs HYPE). ({[}\textbf{c7.45}{]}{[}EXAMPLE and IDEA{]})\\
\item
  A HYPE (hypothetical/hype) is an EXAMPLE in-principle (vs USE-CASE).({[}\textbf{c7.45}{]}{[}EXAMPLE and IDEA{]})\\
\item
  A LINK is a WORD that relates at least two EXAMPLES, IDEAS, or describes ACTIONS, ROLES, and transformations. ({[}\textbf{c7.46}{]}{[}EXAMPLE and IDEA{]})\\
\item
  An IDEA is a set of criteria that LINKs EXAMPLEs as similar (vs not). ({[}\textbf{c7.46}{]}{[}EXAMPLE and IDEA{]})\\
\item
  A PLAN is a BET on a winning arrangement of WORDS that result in the GOAL. ({[}\textbf{c7.23}{]}{[}WORD as BET{]})\\
\item
  A SCIENTIST BETs new CONTENT against the UNKNOWN, evaluated for their net gain for NORMs. ({[}\textbf{c7.47}{]}{[}ROLE Scientist{]})\\
\item
  A RESEARCHER is a SCIENTIST WRITER, a data-collector and hypothesis tester. ({[}\textbf{c7.48}{]}{[}ROLE Scientist{]})\\
\item
  A SCHOLAR is a SCIENTIST READER, curating toward theory development and COMM. ({[}\textbf{c7.49}{]}{[}ROLE Scientist{]})\\
\item
  A LESSON is an ordered set of prompts, usually three to four, up to 60 minutes.(\hyperref[intro-lessons]{\textbf{c8.19}})\\
\item
  A NOTE is a general PLAN for REVISION. (Available on request.) ({[}\textbf{c8.33}{]}{[}NOTES and PEERs{]})\\
\item
  A RECIPE is a DOC format that PRIs CONTENT (IDEAS and LINKS) over STYLE. ({[}\textbf{c8.37}{]}{[}DOC types{]})\\
\item
  COMM is the exchange of WORDS between PPL. (\hyperref[style]{\textbf{c9.45}})\\
\item
  VERSIONs are variations of STYLE on a identical CONTENT, implicitly in temporal order. ({[}\textbf{c9.61}{]}{[}PRIs differ{]})\\
\item
  A PLAN for a DOC is an earlier VERSION of the (same) final DOC. ({[}\textbf{c9.62}{]}{[}PRIs differ{]})\\
\item
  PITCH is DOC STYLE for why to READ. ({[}\textbf{c9.65}{]}{[}NORMs STYLE{]})
\end{itemize}

\section{Draft}\label{draft}

This is an early copy of Skilled Reflection, and much of it is still in development. Thank you for your patience. I welcome any and all constructive feedback (\href{mailto:skilledreflection@gmail.com}{\nolinkurl{skilledreflection@gmail.com}}). Check back later for a revised and updated version.

\section{Version}\label{version}

When I first set out to write SR as a book I quickly realized a problem could hang me up forever, both simple and superficial.
It is the problem of what order to present the information, so that its truth is easily learnable and most evident.
In Version 1 of this book, I deliberately set aside this problem, in order to name and contain the system of thinking.
Version 2 takes the first step in remediation by streamlining all essential terms and links of the system into a shared frame, called ``the Garden'' (introduced in version one). Even still, in this respect version 2 is only a working prototype at best.
I pen this Forward at the advent of betting my life on Version 3 (2022-08-30). I cannot promise if / when it will surface, but if you are curious, do not wait. You can find Version 3 already here, between the lines and at their intersections.

\section{Dear}\label{dear}

Dear, read this book because you're ready to.
You are going to die someday. This book is for you if there is something in the meantime worth maximizing.
Please engage your life. This book is meant to help you understand and take action, starting right now.

\section{Test}\label{test}

Please get a pencil and note the date and time, right now: \_\_\_\_\_\_\_\_\_\_\_\_.
The words in this book are meaningful only when you apply them to your life.
If you wrote the date and time as requested, you have proven enough to yourself that you are open to trying new ideas. you are prepared to grow and continue reading.
If you left the line blank, you might tell yourself you are open to new ideas. You will be right when you complete the task.
Either write the time above or stop reading, and take a better action to improve your life, right now.

\section{Game}\label{game}

From now on, every word you (or i) say is a plan for action.
That plan is meaningless until you take that action, whether good or bad.
Therefore, spend less time worrying about why this word or that plan.
Spend more time deciding which will be the next moves in your game of life.\\
Sincerely, An0ther

\section{Camp}\label{camp}

This book presents plain and obvious facts of life as the antidote to modern complexities of life.
In most cases, the relevance can be seen more clearly if you \hyperref[words]{simulate} the experience being described, or explicitly consider it in your life (see \hyperref[read]{role of reading}). Let's practice. For example, consider going on a camping trip.

\begin{verbatim}
1. Look at the clock and determine when you should expect to arrive.
2. Whatever that time, when you arrive, you will likely be over- prepared, under-prepared, or 1-3 hours late.
3. or you might be an experienced camper.
\end{verbatim}

Either way, this is your only shot at a lifelong camping trip for yourself on earth.
My guess is you are very likely late to your own life.
The games we (people, scientists, etc.) play to explain away this fact, end with this book.

\section{Pitch}\label{intro-pitch}

To understand this book will require you to find and repair the disconnect between your mind's image of time, and the one you are physically bound to.
If you fail, you will die in two worlds, the world you imagine you could have lived, and the one you did.
You, however, would like to ace your own life.
Part one of this book (the garden) gives you all the general answers, and the basic tools for finding the answers particular to you.
Part two teaches you how to (re)build your tools.

\section{Eat}\label{eat}

The one thing up to you that this book will not do is draw the line between need and want.
In this book, food refers to needs and cake refers to wants.
You must practice discerning which is which for you, until it is second-nature.
Then, all decisions can be made as follows:

\begin{verbatim}
1. If hungry, eat food.
2. If full, prepare and pack food.
3. If time, eat or prepare cake.
\end{verbatim}

Beware of thinking cake is food.

\section{Homeless}\label{homeless}

Do, don't think. likely at this point you have slipped into thinking this is a book discussing a philosophy, but it is not.
The difference between food and cake is not understood and solved as a mental exercise in your living room, but by experiencing and paying attention.
What's more, if you don't have a living room, cake can be painfully obvious. be homeless.

\section{Food}\label{food}

The odds are you aren't homeless and will avoid becoming so.
As with any hypothetical, practice finding the nearest analogous, personally meaningful experience:

\begin{verbatim}
1. You couldn't find your car in the middle of winter.
2. You went camping without waterproof gear and it rained.
3. For weeks your shower only produced ice cold water.
4. You were isolated in your house for most of two years during a global pandemic.  
\end{verbatim}

Whatever the circumstance, your habit protests for a change in plans.
However small or profound the suffering or its threat:

\begin{verbatim}
1. There is initially terror, frustration and suffering.
2. Enduring, what you consider suffering will change.
3. You will need far less food than you thought, giving more time and quality cake you really desire.   
\end{verbatim}

\section{Plan}\label{intro-plan}

In the first half of this book is a conceptual metaphor described as a garden.
It is a mental model of reality, designed to bring clarity.
Your garden is your goals put on key ideas of this book.
You need the garden to combat the idealistic side of your mind, inclined to ignore time and your self, which sets you up for failure and suffering.
The second half of these chapters elaborates explicitly on the key ideas.

\section{Lessons}\label{intro-lessons}

At the end of the book are a series of \hyperref[zero-or-one]{lessons} for identifying what matters, what's wrong, and what's next in your life.
They are designed to be useful at any time, and especially repeatedly.
Don't wait to finish the book to try one out, because they are pivotal to skilled reflection.
However, don't do too many of them before at least reading the relevant section that explains the format of {[}skilled reflection{]}{[}skilled reflection{]}.

\section{Definitions}\label{definitions}

Pay attention to how words are defined, as it often differs slightly from conventional use.
For example, ``work'' is an idea referring to ``careers'', but not `yard work'.
These nuances are the result of careful decisions about the food and cake of the ideas.
You can find a list of the key ideas in the {[}index of definitions{]}{[}index of definitions{]}.
Each chapter deals with a distinct perspective or dimension of life and the goals related to it. what is said is what can be simplified for generalizable guidance.
Like food vs cake, the habits relevant to your goals can be closer to food, or excessive like cake.
In the context of definitions, they will be referred to instead, as ``good'' and ``bad'' (for example, {[}c9.60{]}{[}pris differ{]}).

\chapter*{The Garden}\label{the-garden}
\addcontentsline{toc}{chapter}{The Garden}

\chapter{Self}\label{self}

\section{Garden}\label{garden}

Imagine your life takes place in a garden.
Visualizing your life as a garden allows you to ground actions and plans.
The mind analyzes well with visual space.
In this garden we'll consider all your biological needs -- eating, sleeping especially -- as literal.
Your plants are goals, and gardening is your plan.
Food and cake exists in the form of fruit from plants.
Whatever it takes to obtain your cake you can consider as an act of gardening.
Every one of your life ambitions, goals, is buried in the ground as seeds.
Finding a life partner, acquiring a car or promotion, or resolving a conflict with your childhood friend are all matters spread across the dirt in need of water, sunlight, and--most important--your time.
Other goals, like your favorite hobbies, or tv shows, take little work and grow almost anywhere.
Whether they are dear to you or weeds, if they take your time, they are plants in your garden.
Except that they start as seeds.
And what do you do with seeds to make plants? You garden, dear human, you garden.
All actions toward a plant are according to plan.
All actions toward your self are to be considered literal in the garden metaphor.
For example, if you need to eat, or go to the bathroom, you physically go to the place and literally do the deed.
(There will be no pooping on plants, or a pooping plant in this garden.)
Thus, in the garden you are still a normal human bound to a normal space and time, capable of: acting according to bodily need, moving about space, and gardening.

\section{Today}\label{today}

Today, you are looking out across the entire garden.
What has already been planted? What will you add and how much time will you devote to those goals?
Every seed you plant -- even an abstract one like passing a course or making a friend-- implies things in and out of your control like work and weather, or beyond your awareness, like what your boss loses sleep over.
Meditate on the final products of your garden.
Is there a more important plant that deserves your time today?
Meditate on the priorities of your work. do two plants grow better when together?
Make a quick and accurate assessment, but don't sit here all morning thinking about it.
A garden is for working in, to get the food and cake you need.
What do you need to cultivate, today?
What do you want to cultivate in this life?
First, you want what your body and mind best respond to. many of these seeds will be obvious.
It is up to you how curious you are for what better life is possible.
The day is too short to try everything, and many things are not worth trying.
There is not always harm to planting something that does not grow, but sometimes there is.
You could run out of space, time, or energy that could have been devoted to better crops;
You might get a weak yield for trying to grow something not suitable in your climate, or
Find you have no taste or stomach for the fruit.
In the garden, the costs of goals are grounded in their most basic sense: time and attention.

\section{Goal}\label{goal}

The following are core concepts for application beyond the garden.
A goal names an ideal future reality, comprises a plan to increase its likelihood, associated with chud.
These three things encompass all cognitive activity related to reality.
Call anything that could be true, but isn't right now, a goal, and any effort in support of making it true a plan.
It is \emph{your} goal if you aim to increase its probability of being true.
A goal can be finite or persistent, singular or recurring, or realistic or not.

\section{Plan}\label{plan}

If you will need to do something before it is in your hand, it requires a plan.
A plan is any combination of words where the writing or reading of them aims to improve the chance or decrease the cost of achieving your goal.
Most plans are instructions (instr) for learning, making, changing, reminding. product instr are a plan to maximize satisfaction and minimize liability. for example, lamp instr will likely suggest not to put it in water.
``A doc is a plan for goals to complete a linguistic product . example docs include a resume or epitaph on a tombstone.''
``Lessons are a plan template for thinking. . read and complete lesson steps in order.''
Tools are physical or mental objects that help you accomplish your goal.
What is materially required--muscle energy, cash, brain power--is cost.
All costs can be reduced to time.
Cost is any action that has a direct exchange.
Costs often are secondary objectives within a plan.
For example, you might buy a tool (cost).

\section{Self}\label{self-self}

Every day, as a human, you spend time and attention just meeting basic needs: eating, drinking, pooping, and sleeping.
These actions are subgoals of maintaining self, the category of basic human goals.
Food is the necessary and sufficient material to satisfy these needs.
Cake is anything (else) fulfilling, that you are inclined to pursue, including the reason for living.
The goals of self are what's necessary and sufficient to obtain and manage food and maximize bet outcomes (\hyperref[bet]{c5.28}). mind, healthy relationships.
Separate food from cake.
Be careful not to mislabel cake as food. cake may appear sufficient but ask yourself if it is necessary.
Build routines for health and fitness, physical safety (shelter), a sound
Master the self (survive well).

\section{Home}\label{home}

Home is the collection of materials and tools needed for self.
Within a home related actions are grouped into rooms.
For example, a kitchen, is for eating, storing, preparing and disposing of food.
Home also includes all physical material and resources you own, including money, and environment and basis for habits.
A good home is clean and aligned for food.
A bad home facilitates bad habits and takes more time to maintain.
All your life can be very well characterized by a series of purely numerical data:
Event timestamps associated with a location.
The average movement over time, consider as habits.
The average location is home.

\chapter{Priorities}\label{priorities}

\section{Calling}\label{calling}

Your calling--your highest priority--is yours to name.
If you overcome your doubts and see reality plainly, you will want peace.
But we all have doubts, so you want some form of dignity and affirmation.
Tell me your calling, student, and subtracting from peace, i'll tell you your doubts.
On mastering my book your life will be regimented to afford your calling.

Tell me the missing chapter that gives peace instead of regimen, so that i may find mine.

\section{Cakes}\label{cakes}

In your garden you have the seeds and sprouts of the best you could want, and the most important of what you need.
There are two important things to keep in mind, however clear things may seem right now.
The first is to always leave room for your goals to change, if and when you achieve them. think of what you can grow from the seeds of the best apple in your crop.
The second is that much of what you could achieve, you won't, before the day is over, nor before you die.
Before there was a garden, you lived as if you wouldn't die. at least some of your goals assume \emph{tomorrow} is a good time to start.
I did, at least. the list of my goals grew like a wild vine and i found myself overwhelmed by what else i wanted, and yet an inability to equalize various
The very difficult task of your priorities, at their core, is to establish the optimal intersection of finite time and infinite, ever-changing desire.
For example, sleep. about 4 hours daily is a vital priority for everyone (food), but sleeping in on the weekends is usually not (cake).
In your garden, you are limited in what you can do in one day, and in one lifetime.
Even if you build machines to work fields that stretch out to the horizons,
You can only eat so much in a day and a lifetime.
How do you most want to spend your time?
If your goal is to obtain something you do not yet have and have never tasted, be careful. it may not satisfy you how you imagine.
Suppose your goal takes 300 hours of labor.
You might spend 10 years accruing the time on task, or immerse yourself and achieve your goal in under 1 year.
Any true goal has a specific optimal deadline.
Today, you grow the best apple you can,
Yet having tasted it, you might discover a more refined goal is closer to your heart.
Whether this apple is food or cake, the more you align with yourself, the more you will see your goals differently than you first conceived of them.
Leave time in your life to achieve the next goal.

\section{Death}\label{death}

Many ppl believe in seeds with special powers, like true love, or winning the lottery.
One seed grows the duration of your life, which you rarely notice but are often aware of.
When you eat it, the sun sets for the last time and you bite the dust. it happens to all of us.
Think about this, but don't dwell on it.

\section{Sick}\label{sick}

Maybe you are sick, down or depressed.
Then today, do not look out onto your field as a day to work.
Today your field gives to you, and you should take what you need. this is what it is for.
Do not beat yourself up for what you cannot give today. this assumes your field is only meaningful for its future.
And before you're well again, remember this moment well, so you may
Prevent the urge to be disappointed the next time you are less than productive.
If you must work to feel worthwhile, work to budget for a fact of being human:
You are always sometimes less than productive, and occasionally your setbacks come in unfortunate forms, at inconvenient times.

\section{Attention}\label{attention}

Think of what's happening in your garden this very moment: nothing, while you're sitting here thinking.
Reflect only enough to improve your actions.
Pri is the optimizing and managing of all your goals.
All goals rely on a single, limited source of attention, and existence: you.

At its best, this limiting fact drastically improves the reality of plans and goals.
At its worst, it is the most common oversight responsible for bloated {[}chud{]}{[}c.h.u.d.{]}.

The goal of attention is to think; to link your reality here and now with the word that best moves you toward your priority.

Goal importance partly depends on context (c3.6).
Group your goals by shared actions, places, or times. make/group/keep,
Align self and cake.
Ask what it would take to change (if possible), and whether change is worth it.
Savings occur anytime you complete an action that serves various goals.
An example is grouping (c3.35.4).
Apply to ideas, plans, goals, or home things.
Time is the constant (or denominator) for forces, forgetting, pri, and lib.
Do not ignore time.
Estimate durations accurately for goal accomplishment.
Study time to learn reality, self, and their link.
Continually assess whether goal benefits outweighs costs of time.
Maybe it seems unnecessary to represent work satisfaction and relationship quality in terms of carrots and tomatoes.
When you get bored, ask yourself why attempting to understand and define yourself bores you.
What is happening in your life, and what is in your control if these are unclear?
If you are not reflecting on your life, your garden is a foggy labyrinth, and you are a drunk gardener wearing oven-mitts.
You are pushed like a sail by any wind, the forces of reality.
A force is any cause of change.
A force underlies every action involved in a goal, yours or otherwise.
There are forces within your control, and forces outside.
Maintenance is the cost of force to neither move toward nor away from a goal.
Alignment is when norm or natural force causes your goal to be more likely. moment
To align is to adjust your goal to be more like another force, usually one acting against your goal.
Study how forces work against you.

\section{Now}\label{now}

Now is the only use-case of shared space, in time.
A moment is any now that is not the latest, either future or past.
A record of now becomes past `nows' of variable use-case.
That is, we might reflect on the previous nows that we shared.
Note that past nows do not maintain the same level of rigor. they are degraded to hype when mis-replayed, misunderstood, or accessed under changed priorities.
In seeing passage, we might speculate on future nows to be shared.
The discourse of past or future is always at least partially hype, but real (use) when bet on.

\chapter{People}\label{people}

\section{Places}\label{places}

``Just like you, others are trying to figure out what goes on in their own garden.* people are different versions of each others' priorities ({[}see also c9.70{]}{[}norms style{]}).''
Amid the goals and forces of this chapter, none of what is described includes using words. reader, you will begin to understand their nature in the next chapter, bet, but as it pertains words that people use with each other, not until the last chapter, comm.
Consider people in terms of the four places of your life:

\begin{verbatim}
1. On your garden dealing with your self.
2. On someone else's garden.
3. At work making something for others for money.
4. Getting something someone else made (market).
\end{verbatim}

Sometimes you are at your garden's edge looking elsewhere.
For example, when you read an email or social media feed, as if looking across your neighbor's plot,
Studying how to bring your lettuce back to life, or celebrating your friend's successful pumpkin patch.

\section{Gate}\label{gate}

Despite garden-based priorities, most of your time is spent outside your garden, and mostly for work (about 80,000 hours in your life).
Acknowledge the people in your priorities, yet don't let them distract.
Notice how long you leave your garden and to where.
Spend only as much time needed away for your priorities.
For example, most time on others' gardens, and any more than the minimum at work or market is cake. when in doubt, go home and stare at your plants. at least that cake is free.

\section{Work}\label{work}

Where do you get seeds for food from and how did you know to plant them?
From other people, right? No silly, you don't know how to garden a sandwich!
Lucky for you, many of the most important crops you want are already grown, prepared and handed to you--in exchange for money.
Money is traded for maintenance of, or insurance for self and cake, like apples, miracle medical procedures, and a toilet to take your poop somewhere else.
Work is performing a specific task on a collective garden, like an institution's in exchange for money. it is a pre-arranged visit to another garden for a specified time.
Whether or not you like it, or it directly fulfills garden needs, work is made to serve norms, not you.
Norms are all actions assumed of (or about) the ``average'' person. they are the web of forces of all actions of all people, including work, religion, popular attitudes, and doubts.

\section{Analysis}\label{analysis}

If you have a job that pays you to think, your mind is hired to make norms products real.
Instead of your reality, your work is to translate norms's reality, chud-adjusted.
Norms assume you will work for money for goods.
Language, agreed usage of words, is made from norms.
Norms push against individuality (except where it provides a lucrative job opportunity).
Relationships are goals about people.
Good ones are aligned with your goals.
Bad ones cost more.
A relationship is a reciprocal pair of bets, yours of them, and vice versa. one is the better and the other is the bet and bet on.
Relationships, especially family members, coordinate many goals for savings.

\section{Others}\label{others}

Maybe they are your friend, and need help with some unruly vines, or maybe you just like their apples.
Sometimes, people will contribute to your goals, as if bringing water to crops you didn't recognize need them. 4. sometimes people will try to water your crops when they don't need watering.
Sometimes you'll return with higher morale or a bag of apples.
Sometimes you'll need to water some crops that need to be watered, because the sight of their neglect cannot be ignored.
The challenges that overwhelm you are the same for others, just not always the same amount nor at the same time.
The place, duration and impact on your priorities are the basic measurements of a relationship.
Pick work and relationships, including friends, that maximize your other priorities including possibly one that maximizes time and money to apply to other priorities.
Maximally align with norms with least compromise to priorities. get along with people.
people are the part of priorities that require the most care, and fewest words.
A good relationship is a contract of the reciprocated actions (use) between two people.
A bad relationship is hyped mutual bets (unaligned with actions), that waste time.

\chapter{Bet}\label{bet}

\section{Reality}\label{reality}

Life is choices over time.
Each commitment is a bet.
Are you moving toward your priorities?
Will you arrive before you die?

\section{Bet}\label{bet-bet}

A good bet aligns reality and priorities.

To bet, name:

\begin{enumerate}
\def\labelenumi{\arabic{enumi}.}
\tightlist
\item
  A step toward your priorities.
\item
  Estimate the CHUD (see below)
\item
  Set a deadline.
\item
  Take the step.
\end{enumerate}

\section{Win-Lose}\label{win-lose}

You win if:

\begin{enumerate}
\def\labelenumi{\arabic{enumi}.}
\tightlist
\item
  You achieve your step before your deadline.
\item
  You are closer to or have more of your priorities.
\end{enumerate}

Otherwise, you lost time.

In both cases, study to improve your next bet.

\begin{enumerate}
\def\labelenumi{\arabic{enumi}.}
\tightlist
\item
  Compare your action to your priorities.
\item
  Compare your estimated and actual CHUD.
\item
  Revise your priorities.
\item
  Start the next bet.
\end{enumerate}

\section{Cost}\label{cost}

Everything in life costs time or effort.
You must be in a certain place, take action, and often give up something.
For example, physical effort or money.

Self is the basic cost of living your life.
Costs against your priorities are the only bad ones.

Costs only pertain to what actually happens.
To bet requires estimating costs.

\section{Habits}\label{habits}

Habit is the basic pattern of actions you have taken,
whether toward your priorities or toward suffering.
Habit is the rational, descriptive summary of your history.
Habit serves as the rational prediction of your bet outcomes.
Regardless of how expensive your habit, a change of habit is always an added and difficult cost.
That is, unless and until the change becomes a habit.

Suppose in the last 3 days I ate at 9 am.
If tomorrow is a day like the last three, habit predicts I will eat before 9 am.
It also predicts that this is the lowest cost action.

\section{Unknowns}\label{unknowns}

Unknowns are all that are fully out of your control (i.e., statistical error).
For example, planning a drive that doesn't account for a train crossing delay.

Unknowns are the difference between your real and estimated costs and habits.
Bet to learn them.

\section{Doubts}\label{doubts}

Imagine achieving all your priorities, right now.
When considering the costs, note to the fears that surface.
Never mind the world, how are you in your own way?
These are doubts.

Doubts are reliable patterns of failure.
See your doubts with curiosity.

For example:
- self-fulfilled prophesies
- over- or under-estimating habits
- anxiety or insecurity
- distrust in your past

A strong bet is partly made in terms of what you think holds you back most.

\chapter*{The I.B.C.}\label{the-i.b.c.}
\addcontentsline{toc}{chapter}{The I.B.C.}

\chapter{Education}\label{education}

\section{Classroom}\label{classroom}

\textbf{Definition:} A classroom is a special case of market and \hyperref[work]{work}; the direct and shared exchange of \hyperref[analysis]{mental work}, for the purpose of improving individual \hyperref[priorities]{priorities}.

No one chose to be born.
Everyone begins life with their own unnamed and unanswered problems.
Books are an author's answers.
They can only tell a student what the answer is not (quite).
Instead of teaching how to read between the lines, let the student define the problem through their reality and life goals.
Let them write the plan and teach them only what is needed to succeed.
Imagine a perfect course exists, designed to teach you to fulfill your specific ambitions.
Every aspect of what you need to know, that is known and communicable, is the only thing written.
Everything that cannot be known but must be discovered or practiced, is laid out as a set of instructions, described in the words that maximize the learning opportunity, and your progress.
Rather than a course in a classroom, the perfect class is a manual to reference as you live your life, or at least until you've internalized its contents:
When to take a break to strategize your decisions, lessons on what opportunities to watch for and resist, and so on.
In this book, any social role you wish to take on, artist, engineer, therapist, insurance salesperson, reliable partner, would be customized intimately, curated perfectly for what you need.
Any relationship or interpersonal skill that is realistically possible for you is preceded with the guidance and education that prepares you emotionally to choose the right experiences that set you up to be most likely to find and make the most of opportunities to share yourself with another.

\section{Prompt}\label{prompt}

Is there a better version of you?
Are you capable of moving toward it?
Are you ready?
If your answers are yes, then you are a student. to live a better life, you will commit to change your actions.
First, you will learn what changes are needed.

\section{Admission}\label{admission}

This is a boot camp whose purpose is

\begin{itemize}
\tightlist
\item
  To train adults to think harder, clearer and more effectively.\\
\item
  To produce intelligent solutions for personal and social puzzles.\\
\item
  To have a higher cognitive discipline.\\
\item
  To instill shared cake about reason, thinking and discourse.\\
\item
  To empower.
\end{itemize}

It is a training program designed to break down bad habits of thought, and build good ones while immersed here, a culture of rational thinking isolated from the outside world.

\section{Challenge}\label{challenge}

Here you will relentlessly confront your ideas about the world, with instructors drilling intellectual skills (writing essays, arguing for ideas, and developing proposals for action).
You will be trained to move toward your goals with focus, even in the face of perceptual, physical, or emotional distractions.
Over time, students could expect to cultivate a sharper focus on cognitive objectives, resilience to distractions and challenges.
If you graduate, it will be with the ability to identify, develop, and communicate ideal critical, rational arguments, positions or plans (orally or written) given the available knowledge, finite time, and resources at hand (reference material, teamwork).
You will also learn a code of behavior for being a community leader,
Collaborating or competing with an irrational world.
The central requirement for applicants will be a commitment to better understand the self and world.

\section{Student}\label{student}

Students are people with ideas from experience toward selfish goals.
Lessons depend on students' prior knowledge.
A student
Has a goal that can be better named and planned.
Requires time away from betting.
Admits not knowing but capable.
Students learn to represent their knowledge in words; to prefer better, alternative words. a student sees the impersonal as more reliable both for selfish goals and social ones.
A bad student studies to avoid action, or for its own sake.
A student writes ideas and goals to their instructor.
Your goal is clear. you want:
The strength needed to take the right steps and make a habit of
Keen eyes to estimate the destination and correct course, and
A focused mind to steady the foot.

\section{Consent}\label{consent}

This is an in-person immersive experience.
You only really learn what you need to know to be who you really want. you won't learn unless you cannot escape needing it; in a dedicated environment that fosters acquisition, minimizes interference.
Whether you're here for the 7- or 30-day experience, you will work hard every minute.
For every minute of lesson on my time, students are to provide two minutes of writing, either toward others' learning or in direct application toward their goal.

\section{Course}\label{course}

Dedicate to identity growth. Be:

\begin{verbatim}
1. Quiet except when tasks require verbal response.
2. Receptive to work and feedback provided by the instructor.
3. Committed to producing genuinely inspired ideas, working.
4. Respectful that all are equal in voice and ear.
5. Receptive and responsive to prompts and observations.  
6. Motivated to describe solutions that benefit others.  

Learn:  
    1. Precisely and only what is needed to trust a clear picture of  
    2. First, how to make a plan, a map of who you are, and who you  
    3. Repeat part one of this book, until your plan is good enough to act on. 
\end{verbatim}

This will be the beginning of change, and the first test of your commitment.
Your performance is evaluated simply: whether or not you
end up eating, sleeping, thinking, talking, and acting differently.

\section{Think}\label{think}

Think, student.
Do not take notes, simply pay attention.
Everything i say is meant plainly.
If you get confused, forget it, and pay attention to right now.
Our goal here is thinking. thinking happens in your heads.
Right now your job is to think about the truth you see in what i say.
A student has two actions to think: read and write.

\section{Read}\label{read}

Read, student.
A reader is a listener and observer.
Reading is the same as listening to an instructor,
Except that the speaking pace doesn't determine how fast you have to think, or remind you to
Pay attention!
Your attention cannot be trusted on its own,
So lose distractions
Like your smart phone.
You will befriend the simplest scientific instrument, a clock.
The clock is a cue to think.
When it goes off, get back on task.
The clock will babysit your unreliable attention.

\section{Truth}\label{truth}

Your goal in reading is to isolate the truth from the lie.
Try reading this sentence:
``Everyone is best off running weekly until they die.''
You're thinking, ``this can't be true for everyone, so it's a
Not so fast.
There are many components to this idea, and likely many that you
Often a lie becomes true just by changing the pronouns in the
Consider this revision:''i am best off running weekly until i
Perhaps now the words are more truthful to you.
Doing this makes a clear relationship between your belief and the author's.
Becoming smart is the discipline of understanding how you relate to others.
When reading, dismiss only what you fully believe is an intentional lie.
More generally, read to assess your bet on the words reflecting truth.
Whether to a single word, a line, chapter, or book.
Assign weights (0 to 9) to what you read, to complete the following prompt.
I bet this is true for:\\
0 = not even the author.\\
1 = only the author.\\
2 = the author, me, and a few others.\\
3 = us and 30\% of everyone else.\\
9 = us and about 90\% of everyone else.\\
For every assertion and assumption you read, begin assuming it is a ``9'', working backwards according to evidence you hold.

\section{Lesson}\label{lesson}

The following mini lesson illustrates the risk at hand--involuntary comprehension, and the benefit at stake from deliberate reading.
Consider the following quote, ``change your thoughts to change your life.''
Comprehension is involuntary.
You cannot help but recognize meaning when you read.
This means that you likely thought the quote was largely
Reading as proposed here, is strongly voluntary.
For example, reread the quote, this time assuming it is reasonable, serious, and valuable.
The statement is the core assumption to this book, and any psychological theory.
For example, freud's talk therapy was, in his time, the radical idea that words could fix ppl.
Therefore, do not waste the opportunity to consider a truth that could change your life by failing to entertain a simple assumption.

\section{Say}\label{say}

How will your life look if you put it on paper? Like a bunch of words.
How do you change it? By deleting the words with lies, and replacing them with better words.
The right words will change your actions and your life. to live a better life, starts with your words.
In order to do something about thoughts, we need to think on paper. {[}you'll write a lot. you'll delete a lot. you'll get good at writing.{]}

\section{Write}\label{write}

Bet on words. a writer invests time and energy to map feelings onto words.
Good writing is discovering, curating, and applying insight. bad writing has an author; ad hominem.
Revisions also make you a writer. when you revise words (yours or others') to maximize your bet, you are a writer.
Separate thought and self (author), by betting explicitly.
State your assumptions, do not justify them. -is-style-bad
Replace ``I am'' (i.e.~authorship) with tag words to denote your audience (e.g., veteran).
Strive for comm-content and brevity; write only valuable bets, or words that manifest valuable bets. prioritize understanding over original writing.
Writing for \hyperref[library]{people}, relatable, depersonalized, objective words, minimizes rot, maximizes plan utility.

\section{Instructor}\label{instructor}

I am an instructor, a guardian of true ideas and writer of a general plan (this book).
An instructor reads, and enforces bets on links toward a plan.
My goals are to
Minimize student effort and time to write.
Read for cognitive biases, illogical appeals, and imprecise
Reward arguments based on (truth:) reality, self, and cake.
Reward ideas shared (vs kept).

\section{Answer}\label{answer}

Which lines from this book would you {[}bet on{]}{[}how to read{]}, written as-is, or revised for you?
Assign 100 dollars (total) to those lines according to their relative impact, and you are an ibc student; an author of shared truth.
Suppose a representative set of ppl did the same thing. here's a game to help imagine.
Each student's bets go to a general pool.
When bets on a line reach a critical mass, the pool of sub-par bets is split among winners.
Cash is divided.
Also, points are recorded, to incentivize a deeper purpose:
A game, where the prize is
The collective revision and authorship of this book.
After x wins, your name appears on the title.
High scores are on the acknowledgements page.
This is the process and meaning of reader and writer; bound directly to material and action, irl.
Whereas the exchange of money incentivizes the revision of truth, the content itself would be the premise to the next book, {[}the answerword{]}{[}the answerword{]}.

\chapter{Words}\label{words}

\section{Cognition}\label{cognition}

Words describe the world and its conditions.
Each word is a bet about how the world is or can be,
and by extension, the use of words can be improved.
Better words means:
- more practical understanding and expectations.
- better bets.
- healthier communication and relationships.

\section{Bet}\label{word-bet}

A word is a bet.
Good words advance good bets.

Words cost almost nothing but a fraction of a second to make or hear.
Reading, thinking, saying, and writing a word perpetuates that word's idea over others, either moving you toward a goal, or your habit.\\
Words said and heard impact both parties' beliefs and behaviors.

Speaking your mind is difficult to do accurately.

Words are good when they prevent or improve actions you would've taken otherwise.
Unfortunately, the norm for unproductive words is tens of thousands per day.
For example, the words in every ad not aligned with your priorities, or thought about a celebrity you will never meet.

\section{Simulation}\label{simulation}

To make good words, understand their unique power.
A word's power is in narrowing attention in infinite ways.
To reliably narrow, a word requires a definition.
A definition is a generalization which depends on examples.
Careful, good words efficiently isolate and simplify alternative choices.
Bad words bring attention to irrelevant generalizations.

A word is good when it leads to better actions.
Therefore, use words to simulate and find such improvements.
You can think through far more than you can act out.

Words can be easily crossed out and revised.

Good words maximize productivity of thought, move you beyond pitfalls of chud, direct attention to pris, and predict reality; improve decisions and make you smarter.

Simulation is the act of iterating between write and read.
Word meanings are not limited to reality, but your imagination.

\section{It}\label{it}

It (i.e., ``reference'') is the idea words attempt to bring attention to.
For example, consider this scenario:

I wanted to sleep until 9:30 A.M.
I also wanted to know how you are doing.
At 9:00 am you call me and begin to tell me how you are doing.
I yawn.
You say, ``I must be boring you.''

Each of these (below) are words to say ``it'', ranked from best to worst:

\begin{verbatim}
1. You're not, please continue.  
2. No, I'm tired.  
3. No, I woke too early.  
4. Don't make assumptions.  
5. I'm not.  
6. I woke at 8:30am, but wanted to sleep until 9:30am, so I think I yawned because I am tired. I can see why you thought I was bored.   
\end{verbatim}

Good words are ones that uses truth to focus and advance communication.
Below are properties that make each choice unique (and less good).

\begin{enumerate}
\def\labelenumi{\arabic{enumi}.}
\tightlist
\item
  Gives a simple, informative reason to return to communicating.
\item
  Does not add or return to communication priorities.
\item
  May not be relevant or true.
\item
  Focuses on others' mistake and issues a command.
\item
  Begs the question. I.e., ``Then why did you yawn?''
\item
  Focuses on the distraction and makes the same kind of ref which the other did and which caused the confusion.
\end{enumerate}

\section{Write}\label{word-write}

When you have good words to indicate It, do not hold out for better words.
Quickly name the relevant links for It and estimate the loss.
Repair your words only when they pose a problem to the communication priorities.

It, and the actions behind it are more important than who said what (i.e., authorship).
Say the truth or be quiet.
Resist and ignore politeness and impoliteness.
Use words for decisions, not emotions.
Use words to facilitate priorities.
Do not write plans you won't follow.
Words without plans are fantasy and entertainment.
To choose words is to not choose to take action in your life.
Stop talking when action (or listening) is needed.

\section{Read}\label{word-read}

Read when you doubt options available for actions you will take.
Limit the duration or words needed to assess consent. (?)
Distrust ego, and take nothing personally.
Investigate the empirical cake of words.
Ignore style.
Do not let others' words abuse your time or understanding.
Assert boundaries against exaggerated or excessive words.

\section{Example}\label{word-example}

The real world is particular; each experience is an example.
A word describes a set of similar experiences.
1. WORDS are never definite and certain in what they describe of the real world.
2. A WORD's definition is a generalization.
As such, words are
1. less precise than reality.
2. better designed to hypothesize and predict.

All words take their meaning from an average of examples.
It is an example.
An example is one event or object.
Examples can be real or imagined.
A Use-Case is an example taken from a real experience.
A Hypothetical (Hype) is an example in-principle.

\section{Meaning}\label{meaning}

Words direct attention in one of two ways:
Links are words that give relationship between two directions of attention.
Ideas: words that refer to examples,
An idea is a set of criteria that link examples as similar (vs different).
The most basic idea classifies examples as ``It'' or ``not It''.
For example, ``blue'' is simultaneously a specific color and not other colors.

Just as costs are only bad when misaligned with priorities, words are only bad when misaligned with It.

\section{Hype}\label{hype}

Hype vs use-case guides what and if to write. a goal is a named hype. good hype are worthwhile goals, sub-goals, or alt plan actions to consider. don't write hype without a bet.
As all words should be toward pris, there are three general modes of pri words:
H0 = history. habit is to history as action is to doc. good
H1 = plan (pri.txt, proj\_doc.txt)
H2 = doubt especially for bet.

\section{Alternatives}\label{word-alternatives}

Alts are any idea which measurably deviates a plan, and are defined directly with respect to the target idea.
For example, bad is an alt to good, as in ``good-vs-bad''.
Any other valid modes, if at all, are in support of maximal revision in these primary ones.
A plan comprises: 1.the decided action 2.the best bet on it
Alternative actions 4.the goal 5.the relevant chud
Good goals and their plans are supported by uses from your own history (self-h), and next best is norm or others' use or data that generalizes.
For example, a plan for a similar goal has previously succeeded.

\section{Letters}\label{letters}

Mnemonics are letter memory aids.
Reducing a word to a letter increases cognitive efficiency, as long as that letter stands for an idea you will frequently encounter.
These letters illustrate kinds of word usage (not exhaustive):

\texttt{f} define; function\\
\texttt{i} synonym, sense, ``i.e.'', ``as in'',\\
\texttt{e} example, data-point, untested data\\
\texttt{h} history (implied h0)
\texttt{h1} claim, thesis. vs h0 or h2.
\texttt{x} cross-references, xref\\
\texttt{d} data or evidence, summary
\texttt{v} version / revise\\
\texttt{vs} an alternative, sibling of a shared parent category.

\section{Notes}\label{notes}

Good notes clarify pri (content) within bets from secondary ones.
In a doc, par is a grouped set of words roughly equivalent to a complex sentence.
It has a primary subject and predicate, and includes any immediately relevant branches from either.
In practice it is between 1 and 4 clauses.
Good par successfully denotes a link between two ideas, with the following form:
A syntactic tree, where\\
{[}Newline{]} is the right path in a fork,\\
First indent ('' -``) is the left path in a fork,\\
Subsequent newline indents or in-line''--'' are siblings\\
Left to right are siblings\\
Double linebreak ends the local tree.
When word order is a left-to-right walk of a right-branching syntactic tree, sentence can be written in lines of random lengths, and read equally as unambiguously.
Restrict reference and vocabulary to simple- or plain- english to reduce amibiguity.

\section{Transparency}\label{transparency}

Make folder and file names i-goals as transparent as possible imply or reveal its (hidden) members.
Promote or consolidate high frequency items or only-childs.
For example: \texttt{ppl-work\ -\textbackslash{}\textgreater{}\ work\ -\ self-home\ -\textbackslash{}\textgreater{}\ home}
When creating or reusing a word for new applications, this principle should be a factor.

\section{Recipe}\label{recipe}

A recipe is one of the best ways to arrange plans.
Lessons and instructions use a recipe format.
The recipe format highlights the ideas and links of your point and minimizes excessive style.
List key ideas.
Describe actions and transformations (links).
Given an idea, estimate relevance to pris, problems, undeveloped plans, and self-maintenance.
Keep docs short enough that the title and content address only one thing.
Save selectively and delete frequently.

\section{Document}\label{document}

A word is a doc when it is saved with a name, at least once.
Notes are words for communication.

Introduce ideas in unambiguous terms.
Introducing a new topic, provide a succinct, distinctive illustration of the point or merit, \emph{in the verbiage you will most likely understand}.
then describe the link with a taxonomic reference, as follows, where each idea read left to right is a type of or label for the preceding category, implying alternatives at each level.
e-psych-teach-2021-unit4-hw-methods\_report-intro ``for this assignment, consider your grandparents.''

latter ideas qualify the topic, psych. the order of terms from left to right should closely correspond to the order of relevant conditional differences that determine: the reader's priorities.
In the fewest words, help a reader know a topic is irrelevant or truly beneficial?

the most similar and relevant concepts to most distinct, rare, and particular to what is being described.
the left-most word will be either implied or actual chapter headings.

the link is the decision of the sentence that requires the most care.
it is simply a bridge, and as such has a basic and plain function.
named links are actions.
over-spelling the ideas leaves an empty link, e-``do\ldots(the
trash)'' consider what is the manner of spending time, if an
action, to ascribing your link. consider the specific change
being undertaken. some actions: take, give, make, buy, sell,
use/eat, work, write, read.
- notice these are all verbs of transfer.
in instr, the link comes last, and it describes a step in revision.
contextualize with an alternative
a plan should inherit or give definitions of new terms.
new definitions should especially be accompanied by a true use-case, to protect against a false problem.
a pri is the motive for life, and itself is only a named spending of time.
a developed pri is an instruction.
one generalized is a lesson.

\section{Scientist}\label{scientist}

A scientist works to win bets against the unknown. they are a professional writer, evaluated on two metrics:
For their new content.
The net benefit on goal outcomes.
A scientist-researcher is a writer, a data-collector and hypothesis tester.
A scientist-scholar is a reader, curating toward theory development and communication.
More will be said about scientist more broadly. here we focus strictly on the aspect of a scientist which is to develop the description of the world, properly.

\section{Framework}\label{framework}

Framework (frame) is a cluster of definitions.
What metric can be used to compare and test bets on words?
Such a metric requires a general framework for cognition.
This is that general framework.
Just as a child learns skills from ``put it in the box'' to ``put it together,'' and ``solve the problem you creatively set up to solve the impossible,'' so too the highest cognitive function --these days a cooperative one beyond the speed or control of any individual-- deserves treatment of its abilities and applications in order of processing difficulty and utility; utility in contributing to itself however it sees fit, but especially in its allocation of finite resources.

\section{Theory}\label{word-theory}

\[ h1 where p = cognition, good.p = p(wit.goal.sort|pri) \]
Cognitive acts should be engaged in as a pyramid of levels (like food pyramid), number of tasks (width of level).
For example, a base height of 4 (levels) = 4 topics on bottom level,
On 2nd, 2 on 3rd and 1 on 4th. where the bottom level is the
For example, assuming a goal, make an action plan (blitz) constrained by number of tasks by type as follows:
Sort (base task)
Name, def, or measure,
Make (minimum viable)
Test (edit, sim, revise)
Use (reliable, helpful)
Additional tasks must ``wait'' to be addressed until a task of the same level becomes complete and removed from the pyramid, leaving a slot to be replaced.
This is a frame that gives a starting point for assessing how and what optimal capacity and boundaries exist for a human as a cognitive being.
Above is a proposed initial bet, arrived at by sim. here's an example good sim:
Name a pri-scope of irl-events.
For example, tasks may differ in how hard and many times actions
The following properties can be varied by sim. 1-task type
Task-level and task-count
For example, 3 levels starting at 10-count, reducing by 2
The goal of sim is to maximize a bet you take (e.g., vs defining

\chapter{Change}\label{change}

\section{Change}\label{change-1}

Change can be good or bad.
Change is always different.
Good change is revision.
The need for change is dictated by, and specific to priorities.

\section{Reason}\label{reason}

The habit of change is good to maintain.
Sometimes the change needed next is unclear.
Both of these are reasons for this section.
This section develops kinds of change that are often if not always relevant, regardless of priority.

Lines of reasoning

\section{Revise}\label{word-change}

Neutrally, a WORD may be replaced by another. As a process which takes effort, do so to improve, and we will call it REVISION. REVISION is the change of WORDS to improve actions of yourself or others.
REVISION is how we know reflection is happening, described here in VERSIONS (v\#) of a response.
For example, consider the following revisions (v1-3) describing this brief argument:

\begin{verbatim}
1.  1: "Nobody follows doctors' orders."\
    2: "My parents do, religiously."\
    1: "They're the exception."\
    2: "You're not exposed to minorities."  
2.  v1  
    a)  1 makes a false generalization.
    b)  2 illustrates a counterpoint, and 1 gets mad.
    c)  1 dismisses it, and 2 gets mad.  
3.  v2  
    a)  1 generalizes from WORK experience.
    b)  2 argues with parents' experience, and 1 gets mad.
    c)  1 dismisses and 2 get mad.  
4.  v3  
    a)  1 argues outside 2's experience.
    b)  2 uses personal experience, thinking it's impenetrable.
    c)  Failing and hurt, 2 insults.  
\end{verbatim}

The example aims to illustrate revisions which increase the amount of responsibility, control, and preventable future behavior on part of the writer, without much change in WORD count.

\section{Plan}\label{change-plan}

Good revisions better describe what we know to have happened, and predict what will happen (again).
Reality is what determines whether each version is better.
Revise prior prompt responses only to help your current prompt response.
If your goal deals with different assumptions about the truth, change the prompt to whatever gets you to write the most helpful words for your goal.

\begin{verbatim}
1.  For example, a prompt referring to University student experiences can be re-framed for non-University experiences.
2.  A common (problematic) assumption is that you are emotionally ready to be SELF-critical.
3.  Do not change the prompt so you can be lazy.
\end{verbatim}

\section{Think}\label{think-1}

Elsewhere, we will establish variations on ultimate questions.
One variant is the superlative presupposition, which becomes an imperative (plan).
For illustration, all questions presuppose and command, ``think!''
As compliment, all words can be seen as an answer to this question (good or bad).
Supposing just as questions can be revised to their superlatives, we don't bother to ask them, and focus on the superlative of answers.
Thus is an example basis of good revision: to bring the particular apprehension to a universally relevant call to action, responding to the implied imperative of all questions, ``think.''

\section{Plan}\label{change-plan2}

The goal of words is to express ideas of maximum cake tomorrow.
Since few words meet these criteria, start by revising toward fewer words.
Make ideas clear and concrete.
Provide just enough context to remember the basis of the key ideas.
Precision depends on purpose.

\begin{verbatim}
1.  List. Name relevant ideas for goal.
2.  Write a plan, (ordered links).
3.  Doubt. Assert the strongest rebuttal to the plan.
4.  Bet. Improve ideas and links by addressing weakness and clarity.
5.  Read. Wager its cake (e.g., Relative to another plan).
\end{verbatim}

When reflecting, reflect only on ``how can i help my future self?'' And impose time and word limits.
A bad doc is stream-of-consciousness.

\section{Prompts}\label{prompts}

Prompts elicit conflicts (truth) between self and norm, to improve plans for your goals.
Time and word limits work together to encourage a balance between reflecting on truth and describing it.

\section{Limits}\label{limits}

Word limits combat needless words and distracting tangents. a word limit keeps your attention.
A word limit is a proxy for a prompt's complexity.
Try to write the exact number of words.
Good word limits require cutting out unhelpful words, change figurative words to concrete, ideally assertive and falsifiable.
For every plan you make: assert a word limit before writing to be reminded of your initial intentions, and be challenged to express ideas clearly. become skilled at using only the fewest words necessary, to reveal and clarify cake.
If you exceed the limit and there is no end in sight, stop and reassess.

Combined with the word limit, the time limit dictates how precise your words should be.
Use extra time to improve word choice.
For example, given a 50-word limit, 1 minute (1m, 50w) encourages free writing with minimal restrictions on quality of thought, while 4 minutes to produce 50 words encourages more careful selection of words.

\section{Lessons}\label{lessons}

Lessons are an ordered set of prompts, usually three to four, up to 60 minutes and 250 words. a prompt's word count is the number of words to be added to your document.
Lessons target inconsistencies between reality and pris.
They are designed to be revisited and revised repeatedly.
A 0w prompt means revise, but do not increase the word total.
The first prompts in lessons are warm-ups to direct your attention. they ask for names of ideas.
The remaining prompts are for thinking, requiring you to make links between your warm-up ideas. done right, you will face some new truths.

\section{Peer}\label{peer}

Peer revision is a powerful learning tool. get answers from others.
Forget who provides revision and how much.
The peer writer has uncompromised objectivity, and liberty to employ truth, however ``harsh''.

\chapter{Communication}\label{communication}

\section{Pitch}\label{comm-pitch}

Communication happens from one being to another.
The function is to direct attention, coordinate, or share information.
Participation is not always simultaneous.
For example, a traffic sign, rereading your own diary entry, a carefully crafted party invite, or a lease agreement.

Consider you meet someone in the woods, and have only one minute to communicate.
Communicative priorities are the difference between your and others' priorities, especially regarding food and cake.

Say:\\
1. properties of the answer, not knowing it,
2. an answer with the assumption of improving it.
3. common constraints to establish a reliable subjective
experience.
2. what cake defines you (in 250 words), until the self can be measured in its variation from norms.

\section{Style}\label{style}

\section{Communication}\label{comm-comm}

Writing a word creates a static record of a word.
No person is identical with a future or past self, with any other person.
All of these relationships are in part unknown.
The problem of communication is the difference in meaning between reader and writer of the same word.
Even if you wrote the word, your later self may read a different meaning for that word.
Good communication attempts to systematically reconcile these issues.
Bad communication takes advantage of them at the cost of clarity and honesty.
People vary in how they apprehend the world, and therefore they can vary in
1. Precise understanding of meanings.
2. Trust (usually writer more than reader).

To improve communication, study the difference between content and style.

\section{Point}\label{point}

Communication is by default a toward-move between people.
To say a word is to bet on a priority of potential good for the other person.
Each response implies consent that communication is worthwhile for both.
Supposing such a shared priority exists, then the point is rarely less important than manners of conveying it.
Instead, request a separate conversation where such topics are the point of the conversation.

\section{Difference}\label{difference}

Priorities between people (writer and reader) differ.
Difference in priorities alter the content of ideas.
Good style is change in words to minimize change in ideas between reader and writer.
Versions describe identical content with difference in style between them.
A plan for a document is an earlier version of the (same) final document.
For example, you today vs you in five years.

\section{Norms}\label{comm-norm}

communication norms assume reading and writing have no intrinsic goals.
A document's arguments for why to read it are a style called pitch.
Pitch can be accomplished, for example, by

\begin{verbatim}
1.  stroking the ego and intelligence of the reader.
2.  framing attacks as agreeable observations.
\end{verbatim}

The maximum common self-priorities across people are the optimal arguments for pitch. e-food vs.~cake.
Brand is pitch that distorts truth, a form of bad style.
For example, consider the goal of describing the properties of \textbf{apeoplees} with the purpose of selling them:

\begin{verbatim}
1.  CONTENT: **Apeoplees are healthy but sugary.**\
2.  STYLE: **Apeoplees are tasty and nutritious.**\
3.  BRAND: **Apeoplees are healthy.**\
\end{verbatim}

People who do not distinguish food from cake will be persuaded by brand more than pitch.

\section{Save}\label{save}

Recognition of familiar words is easier than their retrieval from memory.
That is, reading is easier than writing.
Save writing that has potential re-use value.
Only save documents that you bet will be useful enough later to save time, overall.

To facilitate retrieval:

\begin{verbatim}
1.  Index (list) documents worth rereading.
2.  Assign a number that indicates its relative importance (abs or
    relative weight).
3.  Add tags and metadata for easier sorting.
4.  Make and revise only for high-priority GOAL(s). Record and study retrieval patterns.
\end{verbatim}

\section{Library}\label{library}

A library is the bet that by saving words to be read again later, future priorities are expedited.
A library is the complete set of documents over your lifetime.
Say fewer words and save reliable priorities.

A good library is a close approximation of self and informative simulation of alternatives.

The risk of a bad library is:

\begin{enumerate}
\def\labelenumi{\arabic{enumi}.}
\tightlist
\item
  writing, saving, and finding documents which mislead instead of help future priorities.
\item
  rewriting an idea that could have been saved and retrieved.
\end{enumerate}

Curating and studying your library affords learning what you don't know you don't know.

\section{Reference}\label{reference}

(Reserved for future development.)

\section{Answerword}\label{answerword}

The complete set of everyone's library describes all that matters in each and all lives.
The efficiently compressed content of this library produces a distribution of variation along a median library.
This is the real point of a library.

This is the end of the book.
Put it in your library.
Now return to your garden and work on your priorities.

\section{Zero}\label{zero}

Also known as Reductio Ad Absurdum.

\begin{enumerate}
\def\labelenumi{\arabic{enumi}.}
\tightlist
\item
  Establish all libraries.
\item
  Study and adopt their implications.
\item
  Transmit culturally.
\item
  The known and lived implication of all words makes words obsolete, one by one, until zero words need said.
\end{enumerate}

\chapter{Self-Guided Lesson}\label{self-guided-lesson}

\section{Expectations}\label{expectations}

This is a subset of the lessons, re-ordered to make for an intense, abridged introduction to the philosophy.

I estimate 2 hours of reading and 5 hours to complete the exercises. As long as you are fully ``present'' for the reading and writing portions, you can expect to do well on it.

In addition to the instructions below, the readings themselves provide guidance on the activities.

\section{Instructions}\label{instructions}

You'll notice that the ``Chapters'' are very short-\/---roughly equal to 1 -- 3 pages of textbook reading, for each, AND in plain language. Do not be deceived; the word choices are very deliberate and should be understood carefully. There is very little style and fluff. Even every example and analogy has an important, central purpose. These concepts are not to be read, but to be mastered.

This activity will require two skills that are ``between the lines''.

First, I'm only assigning parts of the (very short) book, so you will occasionally want to read an unassigned preceding section and/or glossary definitions (in the preface). The second is that, given the ``dense'' style, be careful with reading too quickly just because it is plain English. You might just have to pause to `translate' ideas, and think about why or how something makes sense or is relevant for you. Doing so will make the Lessons much more approachable and beneficial.

However, if you're still confused, email me telling me what you've read and considered along with what you're stuck about and I'll be happy to point you in the right direction!

The lessons are designed to be repeated, so when I'm evaluating it, it matters whether it is your first or second (or 9th) time doing them. Just let me know which version it is, and if you feel necessary, go ahead and share the earlier version(s)---for context. I'll only grade the version you ask me to.
The lessons are to be completed in order, and normally without seeing questions in advance. Take this as a recommendation---I won't penalize you for reading ahead. However, if you do read ahead do not let yourself be discouraged or try to compensate for the intended structure. Two, you absolutely must stick to the time and word limits, and for each lesson, do it all in one sitting. These restrictions are in place to stop you from over-investing in the activities, or getting overwhelmed by them.

For all of the lessons below, I want you to apply them to your actual here-and-now self, not a hypothetical one. You absolutely may choose what aspect of your real life to engage and disclose, but I would rather hear about your struggles with cooking, than a made up/hypothetical struggle with a roommate. This will require some reframing of the literal tasks to fit your circumstance (this is also described in the reading). As one example, ``Student in School'' will talk about midterms and spring break, which you should reframe as this ``make up week'' and ``summer break''. I've put a few more notes like this alongside the outline, below.

One final caveat: Reading ahead might actually be beneficial before doing any of the exercises. I say this because Part 3's readings (its topic is ``how to read'') could make Part 1 and 2's activities and reading more meaningful/clear.

\section{Part one}\label{part-one}

read chapters:
0. Preface (Glossary, here, and big picture of what you're learning and doing with this activity.
1. Introduction
8. Revisions (The section ``Skilled Reflection'' is instructions on how to read/do the Lessons, below).
Turn in lessons by Tuesday, Midnight. The first you may do on your own, and the second one you should record and send a link to the recordings for me to review, along with the saved text file(s).
- Lesson B. ``Define yourself.'' You can skip ``Self-future'', since it might be difficult to follow without more context. The last exercise, ``Self-now'' is one you might want to do twice, or at least note that the time-frame allotted will ask you to move very quickly, and its understandable if the first time through it is a bit rough.
- Lesson E. ``Self-Testing''. This is a general-purpose study technique, based on research from exam 1, so hopefully you'll find it directly helpful. I want you to record your work on this one. The timing is important, and as mentioned above, totally expected (and intended) that it does not result in a polished product. The doing under specific rules is precisely the merit of the activity.

\section{Part two.}\label{part-two.}

Read chapters:
2. SELF
3. Pri.
5. BET.

These first three chapters are meant to give you the conceptual framework that the next two are based on. The conceptual framework is what you will use to ``translate'' your current life experiences, that are relevant. Pay particular attention to this last chapter, BET, and the section at the end called CHUD and DOUBT. What's between the lines here, is that ``Doubt'', beyond its literal meaning, is the term for any / all psychologically atypical or maladaptive experiences. This will be helpful to keep in mind when doing the relevant exercises.
Part Two: Turn in Lessons by Wednesday, midnight:
C. Student in School
D. DOUBT. For this lesson you may want to take its literal framing, or adapt it to be more relevant to yourself as described above and in the ``Revisions'' chapter.
You may choose one of these to be the recorded one, and it's up to you.

\section{Part 3}\label{part-3}

Read these parts of Chapter 6:
6.3 Roles --pay special attention to ``How to Read'' and ``Read PITCH''
6.4 Line number 1 and 2, which describe an activity regarding ``100 dollars''.
Turn in this final activity, by Thursday, midnight:
Read Chapter 7 WORDS and Chapter 9 Communication, and play the ``100 dollars'' game on Chapter 7 and 9. Imagine assigning 100 dollars to the most important numbered lines in Chapter 7 and/or Chapter 9.
That is, find the lines that you think are most valuable for yourself (according to rules and definitions in Roles), and then decide how much it is ``worth'' out of 100 total points. You can make fractions. Pick at least 10 unique lines from at least 4 different sections, but no more than 50 total lines. No time limit for this. For each line that you pick, write a short sentence (up to 10 words) helping me understand your decision, however you think is necessary.

For example, if I was a risky better and trusted my judgment well, I might say
- \$30 for Chapter 7-line 4. So many bad arguments assume people say what they mean.
- \$20 for Chapter 7-line 7. I never regret thinking twice about defining my GOALs.
- And then split 50 among at least 8 other items in different sections.

None of part 3 needs to be recorded, just share the text file.

\appendix


\chapter{Zero or One}\label{zero-or-one}

As a human you are endowed with words.
Too much so. The pervasive impulsivity of putting-to-words means you are often deprived
of un-judged, raw feelings of reality.
For a moment, forget those words.
Let's engage with the manners of communicating
using a vocabulary of 2.

Please answer each question with ``0'' or ``1''.
Do not dwell on any one question.
You may assume anything you like about your words,
but the point of this exercise is to think about what is true without words, and how you might accept (or not)
what another can or cannot know about you without words.

\begin{enumerate}
\def\labelenumi{\arabic{enumi}.}
\tightlist
\item
  Which are you: 0 or 1?
\item
  If the person who knew you best guessed your answer,
  what would they report? 0 or 1?
\item
  What would you tell them you are?
\item
  What would the last dog you encountered say you are?
\end{enumerate}

How do you feel about
5. lunch today?
6. the chair you are sitting in?
7. your fingers?

\begin{enumerate}
\def\labelenumi{\arabic{enumi}.}
\setcounter{enumi}{7}
\item
  Where do you live?
\item
  When were you born?
\item
  What makes this question so simple that you can answer it
  with 0 or 1?
\item
  Is lying, or hoping, or reminiscing 0, or 1?
\item
  Have you been telling the truth to these questions?
\item
  Are you honest with yourself?
\item
  Are you here?
\item
  Do you know?
\item
  How do you know?
\item
  What would it look like if you didn't know?
\item
  What if you \emph{don't know} if you know?
\item
  What are you certain of?
\item
  Is this annoying or interesting?
\item
  Are you ready to use words, yet?
\item
  Imagine you didn't know of language beyond 0 or 1.
  You didn't know a deeper connection with others was possible.
  The chance to respond to the following question is
  your 15 minutes of fame to the universe.
  What is your one bit of say,
  written on the ledger next to your existence, or
  on your tombstone: 0 or 1?
\item
  What's your average score? 0 or 1?
\item
  Has this been insightful?
\end{enumerate}

If you wrote 1 for 24, congratulations.

\chapter{Define Yourself}\label{define-yourself}

\section{Routine and Ideal}\label{routine-and-ideal}

\begin{enumerate}
\def\labelenumi{\arabic{enumi}.}
\setcounter{enumi}{64}
\item
  You are how you act. Therefore, establish a baseline to know what it will take to
  get where you want to be.
\item
  One of the most common mistakes of planning is being overly optimistic in
  predicting outcomes.
\item
  This exercise is a remedy.
\item
  Task 1. 5m, 50w. Take five minutes to produce the WORDS that are most likely
  to accomplish your GOALS for the day. This may include describing the GOALS,
  the PLAN, and/or the DOUBTs.
\item
  Task 2. 5m, 50w. Copy those 50 WORDS and revise them according to the
  following instruction.
  69.1. Replace ``accomplish your GOALS'' with ``do''.
  69.2. In 50 WORDS, what should you (or an all-knowing observer) BET on that
  you will do, today?
  69.3. Regardless of what you assumed in Response one, do not write with the
  intention of `motivating' yourself, but to simply describe your day.
  69.4. If this is too hard, simply evaluate what you did yesterday.
\item
  Task 3. 5m, 0w. Compare your responses. Positive differences (more optimistic)
  in your first response reflect bets you are making against your habits.
\end{enumerate}

\subsection{SELF-750}\label{self-750}

\begin{enumerate}
\def\labelenumi{\arabic{enumi}.}
\setcounter{enumi}{70}
\tightlist
\item
  How accurately can you describe yourself?
\item
  Relate to those who care, including your better SELF, by practicing the right
  WORDS.
\end{enumerate}

\subsubsection{SELF-past}\label{self-past}

\begin{enumerate}
\def\labelenumi{\arabic{enumi}.}
\setcounter{enumi}{72}
\tightlist
\item
  Task 1. 20m, 250w. Look back on your life. What are your CAKE, and where did
  they come from?
\end{enumerate}

\subsubsection{SELF-future}\label{self-future}

\begin{enumerate}
\def\labelenumi{\arabic{enumi}.}
\setcounter{enumi}{73}
\tightlist
\item
  Where are you going? Imagine your life plays out
  74.1. Ideally.
  74.2. Terribly.
  74.3. Typically.
\item
  Consider 1. What could go right or wrong, 2. What it depends on and the
  likelihood. 3. Who you will be if it happens.
\item
  Create a ``choose your adventure'', ordered by likelihood.
\item
  For example, (I will:) 80\% make enough money to be content. 50\% marry, and
\end{enumerate}

\begin{itemize}
\tightlist
\item
  50:50 ends well or in divorce, and 10\% major accident leading to depression.
\end{itemize}

\subsubsection{SELF-now}\label{self-now}

\begin{enumerate}
\def\labelenumi{\arabic{enumi}.}
\setcounter{enumi}{77}
\tightlist
\item
  Task 3. 20m, 250w. Make a small table with three columns.
\item
  In the 1st column, leave enough room in the first column for a short text
  description, and a 1-2 digit number for the 2nd and 3rd column.
\item
  Think about a day and how much TIME you spend on various activities, averaged
  over 365 days in a year.
\item
  on-daily activities such as visiting the dentist should be factored in, accordingly.
\item
  List the activities in the first column.
\item
  List the TIME you spend (in hours), in the second column.
\item
  The total of your second column should be 24.
\item
  Task 3a. 2m, 50w. Choose a GOAL you named in SELF-Future with a probability
  of outcome that you would like to increase.
\item
  What change is needed in your daily schedule to make this increase?
\item
  Estimate the re-allocation of TIME required to achieve your GOAL, in the third
  column.
\item
  For example, if your GOAL is to lose weight, you may need to adjust TIME for
  exercise and eat, and other activities if you ``need more TIME''.
\item
  Task 3b. 3m, 50w. List most likely reasons this PLAN fails, especially reverting to
  old habits.
\item
  Task 3c. 10m, 150w. Add steps to create and protect changes in your schedule.
\item
  Task 3d. 1m, {[}decide-date{]}
\item
  BET on this (Task 3b vs 3c). Set a date on the calendar to assess your
  commitment to this PLAN.
\end{enumerate}

\chapter{Student in School}\label{student-in-school}

\section{Part 1: 10 year plan}\label{part-1-10-year-plan}

This is a lesson in curriculum writing.
Imagine a perfect course exists,
designed to teach you
to fulfill your specific ambitions.

What are your top 4 goals for
the next 10 years?
you could work on any part of yourself,
move toward any goal you want.

\subsection{\texorpdfstring{\textbf{Task 1. 3m, 40w.}}{Task 1. 3m, 40w.}}\label{task-1.-3m-40w.}

list four goals,
Each no longer than 10 words

\subsection{\texorpdfstring{\textbf{Task 2. 3m, 40w.}}{Task 2. 3m, 40w.}}\label{task-2.-3m-40w.}

Your current program requires things your 10 year plan does not.
Courses impose a number of
conditions on your life:

\begin{enumerate}
\def\labelenumi{\arabic{enumi}.}
\tightlist
\item
  Fixed time: courses
\item
  Variable time: assignments
\item
  Uncertain time: studying
\end{enumerate}

List four objectives, goals, course or program requirements that are
most on your mind right now.
Whether most frustrating, or most critical
to you, even `impress my professor'.

\subsection{\texorpdfstring{\textbf{Task 3. 5m, 50w.}}{Task 3. 5m, 50w.}}\label{task-3.-5m-50w.}

Spell out the intersection between the two lists.
Think of how each goal in the second list
directly influences (or not) your 10 year
objectives.

Explore the main ways these two lists relate.

For example,
1. If I don't get a B, I won't keep my scholarship, and will have to
get a job which will likely lower all my grades.

\subsection{\texorpdfstring{\textbf{Task 4. 15m, 100w.}}{Task 4. 15m, 100w.}}\label{task-4.-15m-100w.}

What is the least you can do toward the course requirements and still achieve your long term goals?
Consider cutting the work from list two,
that doesn't overlap with one.

Examples:

\begin{enumerate}
\def\labelenumi{\arabic{enumi}.}
\tightlist
\item
  Graphic design is fun but not necessary. I could drop it and
  still graduate.
\item
  If I get a D, I'll still pass, graduate, and be able to start my own
  business.
\end{enumerate}

Add 100 more words to your description.

\subsection{\texorpdfstring{\textbf{Task 5. 15m, 50w.}}{Task 5. 15m, 50w.}}\label{task-5.-15m-50w.}

Negotiate ways to make direct progress on your first list while fulfilling the second list.

For example, my personality class requires a term paper.
To align with my business goal, maybe the professor will let me write
on ``personality traits of Entrepreneurs.''

Add up to 50 words, and revise.

\section{Part 2: Two week plan}\label{part-2-two-week-plan}

\begin{enumerate}
\def\labelenumi{\arabic{enumi}.}
\setcounter{enumi}{50}
\tightlist
\item
  Often you will define experiences by your priorities and immediate
  desires.
\item
  Experiences out of line with your desires are dreaded, and those in
  line with them are eagerly anticipated.
\item
  For example, a college student dreads midterms and anticipates spring
  break. Each are about one weeklong, yet the feeling of TIME passing
  and the motivation to engage with challenges are very different.
\item
  Anticipations can create disappointment, so reduce speculation,
  especially its harmful effects on you.
\item
  Do more of what you like when you most need it, try to appreciate
  what you dread, as it is, including the suffering.
\item
  Task 1. 4m, 50w. List 4 GOALS to accomplish for midterms, before
  spring break
\item
  Task 2. 4m, 50w. List 4 GOALS for the best spring break you can
  imagine.
\item
  Task 3. 15m, 200w. There are only so many hours in a day. What
  aspects of SELF-care can you incorporate during your midterms to
  improve your sanity, and hopefully your midterm performance?
\item
  We know what happens if you don't. You're staring at your phone;
  you're reading random articles on the internet and wasting TIME.
\item
  Let's take that same honest realism and intentionally PLAN better
  alternatives: instead of garbage breaks, how about spring breaks?
\item
  In what strategic ways can you get a head start on your spring break
  GOALS, as a form of SELF-care and therapy during your marathon of
  midterm WORK?
\item
  Write out your midterm GOALS with SELF-care and/or spring break
  strategically built in.
\item
  Task 4. 15m, 200w. Now imagine that spring break more fully. You're
  spending TIME on one of your GOALS, when suddenly you remember
  something, you need to do for midterms.
\item
  Maybe you get an email from your instructor or boss, or there's a
  project that looms large.
\item
  What are you going to do to be okay with this intrusion? Would it help
  to allow some TIME to WORK on it, or is it vital that you embrace your
  needs and resist any urge to do WORK?
  65.1. Regardless of how you might ideally respond, PLAN for SELF-
  compassion for when you might unexpectedly struggle
  emotionally.
  65.2. Write out your spring break GOALS with midterm GOALS
  incorporated.
\end{enumerate}

\chapter{Lesson: DOUBT}\label{lesson-doubt}

\begin{enumerate}
\def\labelenumi{\arabic{enumi}.}
\tightlist
\item
  Imagine receiving the perfect email that meets your needs; addresses
  your fears, frustrations, paralysis. One that gives you the guidance
  only you could dream of.
\end{enumerate}

\section{Task 1. 4m, 40w.}\label{task-1.-4m-40w.}

\begin{enumerate}
\def\labelenumi{\arabic{enumi}.}
\setcounter{enumi}{1}
\tightlist
\item
  List things the email should include feelings, tangible things,
  states, WORDS of wisdom, anything you need to hear.
\end{enumerate}

Task 2. 4m, author name, revise.\\
3. Consider who could write this: a boss, a best friend, a stranger. Add
an author and revise as needed.

\section{Task 3. 15m, 200w.}\label{task-3.-15m-200w.}

\begin{enumerate}
\def\labelenumi{\arabic{enumi}.}
\setcounter{enumi}{3}
\tightlist
\item
  Now write that email. Don't worry about how or why ``they would say
  this.'' Focus on making sure the letter covers the things you NEED to
  hear.
\end{enumerate}

\section{Task 4. 2m, ``your name''.}\label{task-4.-2m-your-name.}

\begin{enumerate}
\def\labelenumi{\arabic{enumi}.}
\setcounter{enumi}{4}
\tightlist
\item
  Plot twist. Russia sends long range missiles that take out U.S.
  infrastructure including the internet. The government issues a full
  lockdown, and you're home alone without Netflix.\\
\item
  Two days in, an email appears on your computer.
\item
  It's the letter! How is this possible? Slowly you realize, you must
  have written this letter. You did!
\item
  Add these WORDS at the bottom of your letter:
\item
  ``Respectfully, (your name)''
\item
  Read what you need to hear, in the true leader's voice: yours.
\end{enumerate}

\section{Task 5. 15m, 0w.}\label{task-5.-15m-0w.}

\begin{enumerate}
\def\labelenumi{\arabic{enumi}.}
\setcounter{enumi}{10}
\tightlist
\item
  Your original author may have had the authority, resources, or
  information you don't. For example, ``Here's 20 thousand dollars to
  build your classroom.''
\item
  Revise to make the letter practical.
\item
  Tip 1: For resources you don't have-- money, TIME, influence--
  accept that indeed, you do not have power over these things, give up
  attempts to control, and welcome your feelings.
\item
  Tip 2: For guiding your behavior--decisions, actions, confidence--
  give permission to take action.
\end{enumerate}

\chapter{``Self-testing''}\label{self-testing}

All of this book is a lesson in simulation,
because this is a book with words in it,
and words are simulations of reality.
This lesson is an exercise to fearlessly
take action. Since you'll be writing,
it works better if the ``action'' to take is,
itself, a writing task.
It could be an upcoming exam, or writing your will.
The exercise works whether you are
75\% done or 5\% prepared to start.

\textbf{Task 1. 3m, 40w.}

Think about the challenging task you've identified.
Finish a sentence like this:
``The purpose of this task is to\ldots{}''
Perhaps ``demonstrate my knowledge of the endocrine system,'' or
``communicate my values and divvy up my valuables.''

So that I don't continue to prime your death, the remainder of this lesson will focus on the youthful experience of learning.

\textbf{Task 2. 2m, 0w.}

Picture someone who doesn't struggle to do this task.
They would score 100\% or write the perfect paper/assignment.
Assume they weren't born able to do this.
How did they get that skill?
Edit the same 4 things to try to make them more concrete and tractable.

\textbf{Task 3. 8m, 100w.}

Assume that ideal performer writes a custom quiz on
the very areas you're struggling with.
Write quiz questions that tests your learning objectives.
For example, I'm working on a paper that deals with
how talking happens in the brain.
A question in my quiz is, ``define the two contrasting theoretical
positions in 15 words, each.''

\textbf{Task 4. 2m, 0w.}

Revise your quiz so that it would take about 30 minutes for your expert to complete, with questions that can be answered in 5 to 25 words.
Add word and time limits for each question.

\textbf{Task 5. 20-25m.}

Take your quiz.

\chapter{Other Pedagogies}\label{other-pedagogies}

\section{Introduction}\label{introduction-1}

What follows is a bibliography to the educational themes in the main text, most of which are described in \hyperref[education]{Education}.

The Skilled Reflection Plan / Intellectual Boot Camp (IBC) intentionally makes very few references to other work. Loosely, here are the stylistic claims, however a more general argument can be taken from \hyperref[words]{Words} and \hyperref[communication]{Communicaiton}.

\begin{enumerate}
\def\labelenumi{\arabic{enumi}.}
\tightlist
\item
  Authorship is an ego-based IDEA that, at best, distracts from the intrinsic right of a reader to adopt and interact with information.\\
\item
  Most definitions and descriptions of something novel (e.g., IBC), if compared to the familiar, are confined to be framed in terms of the familiar. References to the familiar might suggest the IBC \emph{builds on} or \emph{contrasts from} and thus is subject to the same assumptions or purposes of other work / projects.
  Specifically, while IBC is here compared to schools and pedagogies, its core framework is consistent with SR, which aims to empirically (i.e., cognitive science), meaningfully engage all human relationships that rest on information and social capital: therapeutic models, epistemology, phenomenology, ethics and law.
\end{enumerate}

To highlight similarities and differences with past work is not a guaranteed remedy, would be labor intensive, and emphasizes the past.

Nonetheless, there are also benefits.
References provide an approximation of IBC's similarity with reality, which argues for proof-of-concept.
This list excludes mainstream / traditional education models, focusing instead on merits and limitations of alternative approaches to building rational knowledge-bases.

The only requisite assumption of the IBC is that the \emph{student} assumes responsibility (and feasibility) of guiding their learning. All other assumptions are secondary, and subject to empirical testing. This point is re-emphasized at the end of section A. Schools and Pedagogies.

\section{Schools and Pedagogies}\label{schools-and-pedagogies}

Some educational programs prioritize intrinsic motivation and reasoning skills. They are listed here in terms of similarity/value with IBC, followed by differences. These lists are at best illustrative rather than descriptively complete.

\subsection{Fixed content}\label{fixed-content}

\begin{enumerate}
\def\labelenumi{\arabic{enumi}.}
\item
  \href{https://www.sjc.edu/academic-programs}{St.~Johns} (liberal arts college).

  \begin{enumerate}
  \def\labelenumii{\arabic{enumii}.}
  \tightlist
  \item
    Similarity to IBC: The program fosters independent critical thinking and close reading.
  \item
    Difference: A curriculum emphasizing primary texts of great minds. Students read the Great Books''.
  \end{enumerate}
\item
  \href{https://en.wikipedia.org/wiki/Deep_Springs_College}{Deep Springs} (liberal arts college/community)
\end{enumerate}

\begin{itemize}
\item
  fewer than 30 students
\item
  Similarity: academic fundamentals, self-reliance: students work 20hrs a week, and participate in self-governance.
\end{itemize}

\subsection{Quasi ``Open'' content}\label{quasi-open-content}

\begin{enumerate}
\def\labelenumi{\arabic{enumi}.}
\setcounter{enumi}{2}
\tightlist
\item
  \href{https://en.wikipedia.org/wiki/Montessori_education}{Montessori}
\end{enumerate}

Similarity: pedagogy built from intrinsic motivation (curiosity).

\begin{enumerate}
\def\labelenumi{\arabic{enumi}.}
\setcounter{enumi}{3}
\tightlist
\item
  \href{https://en.wikipedia.org/wiki/International_Baccalaureate}{International Baccalaureate} (IB)
\end{enumerate}

``IB learners strive to be: Thinkers, Communicators, Principled, Reflective, Open minded, Balanced, Risk-takers, Inquirers, Caring, and Knowledgeable.''

For comparison using similar themes, an IBC statement in as few words:

Students arrive curious, ambitious. Personal goals sharpen their wits.\\
They skillfully invest in ideas, communication, reflection, and risks.

\subsubsection{about ``Open'' content programs.}\label{about-open-content-programs.}

\begin{enumerate}
\def\labelenumi{\arabic{enumi}.}
\tightlist
\item
  Emphasize critical thinking. (All)
\item
  Sometimes require intrinsic motivation (1,2) and some foster it (3,4).
\item
  Some heavily emphasize reading, requiring a strong capacity for abstraction (1,2).
\item
  To some degree, all emphasize a curriculum not defined by the student. For example, ``core skills'', literature, etc.
\end{enumerate}

Critical thinking and intrinsic motivation are good focal points of an institution. Like (3 and 4), all students have this capacity. What these programs fall short of, to varying degrees, is a strong recognition that the only thing that can be taught is that which the student wants to learn. Any aspect of curricula that is defined by the institution, such as a book to be read, or a concept to be learned, is at risk of failing to harness individual intrinsic motivation. A program that succeeds at teaching a student, is one that begins with the student, before building a curriculum. This is what I call an ``Open'' content curriculum.

\subsection{Value-based programs}\label{value-based-programs}

Some programs are built to guide and frame experience. They may also do some labelling of experience (content), but the pedagogical emphasis is on instilling a value system. These programs tend to offer flexibility for variation in the students' prior experiences and thus motivation, and enable an opportunity for them to adapt their own content. All of these programs offer a value system for deriving content. The basis of these value systems however will vary in terms of it's intrinsic motivation to each individual, and arguably, its empirical basis of deriving values. Strictly speaking, we can say these programs push content much like the first groups. I separate them here, because they are conceptually framed in a way that is very flexible to the changing dynamics of the real world, and the concepts they push are at least somewhat ego-centric.

\href{https://www.landmarkworldwide.com/}{Landmark Forum}

\begin{itemize}
\tightlist
\item
  Emphasizes cognitive bias training, motivational training.
\end{itemize}

\href{https://strozziinstitute.com/}{Strozzi Institute}\\
* Emphasizes self-actualization and self-efficacy

\subsection{Other programs}\label{other-programs}

\emph{Monastic traditions} (religious community)
\emph{Jesuit and/or Buddhist}

\begin{itemize}
\tightlist
\item
  self-sufficiency
\end{itemize}

\href{https://www.selfauthoring.com/}{Self-authoring} by Jordan Peterson
(an online only self-paced writing program)

\begin{itemize}
\tightlist
\item
  Critical self-reflection
\item
  Self-paced, somewhat one-size-fits all.
\end{itemize}

\emph{Coding Bootcamps}/\emph{Referral Programs}\\
These represent the ideal strictly in terms of measurable output, i.e., resume or job-capital. Their downside is cost and inflated self-esteem. Whereas students are taught to trust their ability, they are rarely informed that very little ability is necessary to land a job, which risks mistakened confidence or inflated imposter syndrome once they begin work.

\section{General observations}\label{general-observations}

Here I'll focus on content and value systems. IBC asserts that a value system emerges automatically when a student commits themselves to achieving what they are motivated to achieve.

A ``blank slate'' learning environment is ideal but not realistic. Instead, students begin in a constructed environment that is familiar to their current reality. The instructor facilitates the student's destruction of the environment--either deconstruction of their reality down to a tabula rasa to be re-built again (bootcamp), or to be revised and optimized toward their own desires (skilled reflection).

\section{Ideal Curriculum}\label{ideal-curriculum}

What is a tabula rasa that works for all students? That which is the basis of all problems for all people: the existential facts of self and self-awareness.

All humans must eat, sleep, shit, and die. Students begin with a curriculum that forces them to experiment on these facts of life with the goal of gaining more out of life. If you've read the book, this should sound familiar. The answer to an ideal curriculum begins with Chapter one: \hyperref[self]{The Self}, and proceeds accordingly.

  \bibliography{book.bib,packages.bib}

\end{document}
